\newcommand{\GBMEMRate}{0.500}
\newcommand{\GBMEMWeakRate}{0.999}
\newcommand{\GBMMilRate}{0.991}
\newcommand{\GBMMilWeakRate}{1.000}
\newcommand{\GBMAnalyticWeakRate}{0.9993}
\newcommand{\GBMFitWindow}{$N=64,128,256,512,1024$}
\newcommand{\GBMExactMean}{105.1987}
\newcommand{\GBMExactVariance}{1045.0120}
\newcommand{\GBMExactMeanCI}{[105.0403, 105.3571]}
\newcommand{\NARate}{0.491}
\newcommand{\NAFitWindow}{$N=8,16,32,64$}
\newcommand{\NAExactMean}{105.08796}
\newcommand{\NAExactMeanCI}{[104.89962, 105.27629]}
\newcommand{\NAPayoff}{14.53804}
\newcommand{\NAPayoffCI}{[14.40515, 14.67093]}
\newcommand{\NAWeakInterpretation}{The 95\% intervals include zero at $N=64,128,256$. No weak order is assigned to the complete grid range. The finest-reference payoff is estimated as 14.53804 with interval [14.40515, 14.67093].}
\newcommand{\NAReferenceInterpretation}{The last refinement discrepancy is 0.04016. It is 3.2\% to 8.9\% of the coarse strong differences on the prespecified fitting window. This meets the practical 10\% reference-sensitivity criterion used for that window. For $N=256$, the ratio is 18.0\%, so the finest coarse points are retained as sensitivity evidence, not used for the principal rate claim. The two refinements decrease the strong discrepancy; their paired payoff intervals assess whether the benchmark changes at a detectable weak scale. Neither a small refinement difference nor a zero-containing interval is a rigorous bound on the unknown continuous-time error.}
